\documentclass[11pt,a4paper]{article}
\usepackage[utf8]{inputenc}
\usepackage[spanish]{babel}
\usepackage{amsmath,amssymb,amsthm,amsfonts}
\usepackage[paperheight=27.5cm, paperwidth=21cm,% Set the height and width of the paper
includehead,
nomarginpar,% We don't want any margin paragraphs
textwidth=17cm,% Set \textwidth to 18cm
textheight=18.4cm,% Set \textwidth to 22cm
headheight=46mm,% Set \headheight to 10mm
]{geometry}

% Useful packages
\usepackage[colorlinks=true, allcolors=blue]{hyperref}
\usepackage{indentfirst}
\usepackage{latexsym}
\usepackage{amsmath,amssymb}
\usepackage{enumerate}
\usepackage{float}
\usepackage{multirow}
\usepackage{makecell}
\usepackage{amssymb}
\usepackage[table]{xcolor}
\usepackage{graphicx}
\usepackage{caption}
\usepackage{xspace}
\usepackage{amsmath}
\usepackage{color}
\usepackage{tikz}
\usetikzlibrary{fit,automata,arrows}
\usetikzlibrary{patterns,hobby}


\newcommand{\X}{\cellcolor{gray!55} \checkmark}
\newcommand{\N}{\mathbb{N}}
\newcommand{\R}{\mathbb{R}}
\newcommand{\Z}{\mathbb{Z}}
\newcommand{\Q}{\mathbb{Q}}
\newcommand{\C}{\mathbb{C}}
\newcommand{\fin}{\hfill\vbox{\hrule height 5pt width 5pt }\bigskip}

\newcommand{\NP}{{\sf NP}}

\usepackage{helvet}
\renewcommand{\familydefault}{\sfdefault}

\usepackage{fancyhdr}

\fancypagestyle{plain}{%
  \fancyhf{}%
  \fancyfoot[R]{\thepage}%
  \renewcommand{\headrulewidth}{0.4pt}% Line at the header invisible
  \renewcommand{\footrulewidth}{0.02pt}% Footer line not visible with 0pt

\fancyhead[C]{
\begin{center}
{\Large
\textbf{UNIVERSIDAD DE BUENOS AIRES \\
Secretaría de Ciencia y Técnica}}\\

\bigskip 

\large \textbf{Proyectos de Investigación Básica, Aplicados, de Transferencia e \\ Innovación Tecnológica \\
Programación Científica 2023}
\end{center}}
}

\usepackage{tcolorbox}
\newtcolorbox{mybox}{colback=gray!30,
boxrule=1pt,arc=0pt,boxsep=0pt,left=2pt,right=2pt,leftrule=1pt}

\renewcommand\spanishtablename{Tabla}


\begin{document}

%\title{.....}
%\author{....}
%\date{\today}

%\maketitle

\pagestyle{plain}

\begin{mybox}
\centering \textbf{PLAN DE INVESTIGACIÓN}
\end{mybox}

\medskip

\textbf{Título del plan de investigación del proyecto:} (debe coincidir con el consignado en SIGEVA)\\


\textbf{Nombre del director y codirector/es:} (debe coincidir con el que se consigne en SIGEVA)\\

\textbf{Disciplina/área del proyecto:} (debe coincidir con el que se consigne en SIGEVA) 

\section*{Estado actual del conocimiento sobre el tema} {\small (Desarrolle en 4 carillas como máximo) Mediante las citas/comentarios de las mayores contribuciones en el tema específico publicadas por grupos distintos al propio, debe explicitarse el desarrollo acumulado del conocimiento sobre el tema. Se evalúa la profundidad del conocimiento que el Director (y grupo) tienen acerca del tema que proponen encarar. Concretamente, la originalidad debe referenciarse mediante los interrogantes aún abiertos y relacionados con el proyecto en cada una de estas contribuciones así como con las contribuciones publicadas de los investigadores del grupo del proyecto en el tema específico.\\}

.... \cite{Uriel-tesis,G-J-S-NPC}

\section*{Objetivos e hipótesis de la investigación} {\small (Desarrolle en 4 carillas como máximo) Fundamentalmente, se deberá detallar el problema o situación de referencia en el que se desarrolla el proyecto, o los interrogantes en el campo disciplinario a los que el proyecto se dirige. Se deben enunciar de manera clara las metas concretas a alcanzar en el marco del proyecto indicando hipótesis o postulados o propuestas explicativas de la pregunta en estudio.\\}

...

\section*{Metodología} {(Desarrolle en 4 carillas como máximo) Describir, según corresponda al tipo de proyecto, el diseño experimental, o el procedimiento, para la recolección de información y su procesamiento. Es aconsejable la descripción muy breve de la metodología a usar dentro de cada sección donde se describen las tareas.\\}

....

\section*{Antecedentes en la temática} {\small (Desarrolle en 2 carillas como máximo) Detallar de manera resumida los avances no publicados producidos sobre el tema, tanto individuales como grupales.\\}

....

\section*{Transferencia de resultados} {\small (máximo 2 carillas) Describir el objeto de la transferencia, su importancia, los destinatarios concretos o posibles, y los procedimientos para concretarla.\\}

....


\section*{Cronograma de actividades} {\small Detallar las actividades propuestas con su secuencia o encadenamiento lógico y metodología a usar en cada una de ellas. Consigne sucesivamente cada actividad unitaria.}

\begin{figure}[H]
\makebox[\textwidth]{\parbox{1.5\textwidth}{%
\centering
\begin{tabular}
{ | m{6cm} |  m{4mm} | m{4mm} | m{4mm} | m{4mm} | m{4mm} | m{4mm} | m{4mm} | m{4mm} | m{4mm} | m{4mm} | m{4mm} | m{4mm} |}
 \hline
 \cellcolor{yellow} & \multicolumn{12}{ c| }{\cellcolor{yellow} \bf Meses del primer año}   \\ 
 \cline{2-13} 
 \multirow{-2}{*}{\cellcolor{yellow} \bf Actividad} &  \cellcolor{yellow}  1 & \cellcolor{yellow} 2 & \cellcolor{yellow} 3 & \cellcolor{yellow} 4 & \cellcolor{yellow} 5 & \cellcolor{yellow} 6 & \cellcolor{yellow} 7 & \cellcolor{yellow} 8 & \cellcolor{yellow} 9 & \cellcolor{yellow} 10 & \cellcolor{yellow} 11 & \cellcolor{yellow} 12 \\ \Xhline{2\arrayrulewidth} 
{\small Tema 1.....} 
 & \X & \X & \X & \X & \X & \X & \X & \X & & & & \\ 
 \Xhline{2\arrayrulewidth} 

{\small Tema 2.....} 
 &  &  &  & \X & \X & \X & \X & \X & \X & \X & & \\ 
 \Xhline{2\arrayrulewidth} 

\hline
 \end{tabular}
}}
\end{figure}

\begin{figure}[H]
\makebox[\textwidth]{\parbox{1.5\textwidth}{%
\centering
\begin{tabular}
{ | m{6cm} |  m{4mm} | m{4mm} | m{4mm} | m{4mm} | m{4mm} | m{4mm} | m{4mm} | m{4mm} | m{4mm} | m{4mm} | m{4mm} | m{4mm} |}
 \hline
 \cellcolor{yellow} & \multicolumn{12}{ c| }{\cellcolor{yellow} \bf Meses del segundo año}   \\ 
 \cline{2-13} 
 \multirow{-2}{*}{\cellcolor{yellow} \bf Actividad} &  \cellcolor{yellow}  1 & \cellcolor{yellow} 2 & \cellcolor{yellow} 3 & \cellcolor{yellow} 4 & \cellcolor{yellow} 5 & \cellcolor{yellow} 6 & \cellcolor{yellow} 7 & \cellcolor{yellow} 8 & \cellcolor{yellow} 9 & \cellcolor{yellow} 10 & \cellcolor{yellow} 11 & \cellcolor{yellow} 12 \\ \Xhline{2\arrayrulewidth} 
{\small Tema 1.....} 
 & \X & \X & \X & \X & \X & \X & \X & \X & & & & \\ 
 \Xhline{2\arrayrulewidth} 

{\small Tema 2.....} 
 &  &  &  & \X & \X & \X & \X & \X & \X & \X & & \\ 
 \Xhline{2\arrayrulewidth} 

\hline
 \end{tabular}
}}
\end{figure}

\begin{figure}[H]
\makebox[\textwidth]{\parbox{1.5\textwidth}{%
\centering
\begin{tabular}
{ | m{6cm} |  m{4mm} | m{4mm} | m{4mm} | m{4mm} | m{4mm} | m{4mm} | m{4mm} | m{4mm} | m{4mm} | m{4mm} | m{4mm} | m{4mm} |}
 \hline
 \cellcolor{yellow} & \multicolumn{12}{ c| }{\cellcolor{yellow} \bf Meses del tercer año}   \\ 
 \cline{2-13} 
 \multirow{-2}{*}{\cellcolor{yellow} \bf Actividad} &  \cellcolor{yellow}  1 & \cellcolor{yellow} 2 & \cellcolor{yellow} 3 & \cellcolor{yellow} 4 & \cellcolor{yellow} 5 & \cellcolor{yellow} 6 & \cellcolor{yellow} 7 & \cellcolor{yellow} 8 & \cellcolor{yellow} 9 & \cellcolor{yellow} 10 & \cellcolor{yellow} 11 & \cellcolor{yellow} 12 \\ \Xhline{2\arrayrulewidth} 
{\small Tema 1.....} 
 & \X & \X & \X & \X & \X & \X & \X & \X & & & & \\ 
 \Xhline{2\arrayrulewidth} 

{\small Tema 2.....} 
 &  &  &  & \X & \X & \X & \X & \X & \X & \X & & \\ 
 \Xhline{2\arrayrulewidth} 

\hline
 \end{tabular}
}}
\end{figure}


(Sólo completar el tercer año en caso de solicitar proyecto ``Modalidad 1'')

\section*{Presupuesto total} {\small Realizar una breve justificación para cada rubro desagregado por objeto del gasto (de acuerdo con lo consignado en el ítem Recursos Financieros de la postulación de su proyecto).\\}

....


\renewcommand{\refname}{Bibliografía}
\bibliographystyle{plain-abbr}
%admite tesis de licenciatura como @LICTHESIS
\bibliography{bibliography}

\end{document}